% Build with: latexmk -pdf main.tex
\documentclass[11pt]{article}

\usepackage[T1]{fontenc}
\usepackage{lmodern}
\usepackage[margin=1in]{geometry}
\usepackage{microtype}
\usepackage{amsmath,amssymb,amsthm,mathtools}
\usepackage{booktabs}
\usepackage{enumitem}
\usepackage{xcolor}
\usepackage[colorlinks=true,linkcolor=blue!50!black,citecolor=blue!50!black,urlcolor=blue!50!black]{hyperref}

\theoremstyle{plain}
\newtheorem{theorem}{Theorem}
\theoremstyle{remark}
\newtheorem*{remark}{Remark}

\newcommand{\K}{K}
\newcommand{\E}{E}
\newcommand{\eq}{\mathrm{eq}}
\newcommand{\mle}[1]{\widetilde{#1}}
\newcommand{\hash}{\mathsf{H}}
\newcommand{\concat}{\mathbin\Vert}
\newcommand{\rootnode}{\mathit{root}}
\newcommand{\comb}{\omega}
\newcommand{\idx}{I}
\newcommand{\loc}{\mathrm{loc}}
\newcommand{\code}[1]{\texttt{#1}}

\emergencystretch=1.5em

\title{Zero-knowledge Merkle membership from two evaluation claims}
\author{}
\date{}

\begin{document}
\maketitle

\begin{abstract}
We describe a zero-knowledge proof of one fixed statement, a BLAKE2s preimage together with a Merkle path of depth 32, built for a verifier that pays for every byte it reads and every field operation it performs, and that can hand a multilinear evaluation claim to a precompile. The proof is two sumchecks over a binary field: 98 field elements and one commitment, 2{,}384 bytes. It ends on two plain evaluation claims, one on the prover's commitment and one on a commitment to the circuit, and the verifier never reads the circuit. Zero knowledge is statistical, with error below $2^{-118}$ for an honest verifier. The scheme is implemented in \code{flock::membership}.
\end{abstract}

\section{Setting}

Almost all of a hash-based proof is the opening of its polynomial commitments: Merkle paths, hundreds of kilobytes of them. The rest is small. We therefore assume a verifier of the following shape, and design for it.

\begin{itemize}[nosep]
  \item A \emph{precompile} checks one claim $(\text{commitment}, \text{point}, \text{value})$: that the multilinear polynomial behind the commitment takes that value at that point. The commitment is WHIR~\cite{ACFY25}. Its opening does not ride the transaction: openings travel beside transactions and are aggregated recursively before a block is built, so a claim has a flat price.
  \item Everything else is a contract. It reads the proof from calldata, replays the Fiat-Shamir transcript, and does arithmetic in the proof's field.
\end{itemize}

The task of the protocol is then to reduce a statement to few evaluation claims, in few scalars and few multiplications. Three choices follow, and they are the content of this note: a sumcheck round sends two scalars and the prover sends nothing the verifier can derive; the verifier touches the circuit through three field elements only; and zero knowledge adds five scalars to the proof.

\paragraph{Fields.} $\K = \mathrm{GF}(2^{64})$ and $\E = \mathrm{GF}(2^{192}) = \K[y]/(y^3 + y + 1)$. Witness values are in $\K$; challenges and proof scalars are in $\E$, 24 bytes each. Addition is XOR. We write $\mle f$ for the multilinear extension of a table $f$ on $\{0,1\}^m$, and $\eq(p, x) = \prod_i \bigl(p_i x_i + (1 + p_i)(1 + x_i)\bigr)$.

\section{The statement}

Let $\hash$ be BLAKE2s~\cite{blake2} with 32-byte output. The public input is $(y, \rootnode)$, 32 bytes each. The prover knows $x$ of 32 bytes, siblings $s_1, \dots, s_{32}$ and direction bits $d_1, \dots, d_{32}$ such that
\[
  \hash(x \concat \mathtt{0x11}^{32}) = y, \qquad
  h_0 = \hash(x \concat \mathtt{0x22}^{32}), \qquad
  h_i = \begin{cases} \hash(h_{i-1} \concat s_i) & d_i = 0 \\ \hash(s_i \concat h_{i-1}) & d_i = 1 \end{cases}, \qquad
  h_{32} = \rootnode .
\]
Every input is 64 bytes, one BLAKE2s block, so the statement is 34 compressions. It says that $y$ is a tag of a secret $x$ whose other hash, $h_0$, is a leaf of the tree, without saying which leaf.

\section{The circuit}

The statement is a rank-one constraint system $(Az) \circ (Bz) = Cz$ over $\K$, on $N = 2^{22}$ wires and as many rows: row $\ell$ defines wire $\ell$ as a product of two sums of wires.

\paragraph{One bit per wire.} A free input carries the row $z_\ell = z_\ell \cdot z_\ell$, which forces a bit, and every wire computed from bits by XORs and ANDs is a bit. BLAKE2s is XORs, rotations and additions modulo $2^{32}$; an addition costs 31 rows, one per carry. Because a bit is a whole field element, there is nothing to unpack: the proof needs no ring switching, and a public value is pinned by a linear condition on wires.

\paragraph{Blocks.} A wire index is $(\text{local}, \text{slot}, \text{compression})$, of $11 + 5 + 6$ bits: 64 compressions (34 real, the rest padding), each of 32 blocks of 2048 wires. A compression uses 22 blocks, of three types:
\begin{center}
\begin{tabular}{@{}cll@{}}
\toprule
type & count & computes \\
\midrule
$\mathsf S$ & 1 & the message, $(h, s)$ swapped by the direction bit, and the constant initial state \\
$\mathsf H$ & 20 & half a round: the function $G$ on the four columns of the state \\
$\mathsf F$ & 1 & the output, the final XOR of the two state halves; it also holds $x$ \\
\bottomrule
\end{tabular}
\end{center}
The diagonal half of a round is the column half applied to a state whose rows are rotated, so one block type serves both.

\paragraph{The matrices are local plus wiring.} Within a block, a row reads wires of its own block: that is a \emph{local matrix} $A_T, B_T \in \{0,1\}^{2048 \times 2048}$ per type $T$, the same in every block of the type. The only other rows are copies, $z_\ell = z_{\mathrm{src}(\ell)} \cdot 1$, by which a block reads the previous block's state, a message word (through the BLAKE2s schedule), or the previous compression's hash. The map $\mathrm{src}$ is a product of small maps on the coordinates (bit, word, slot, compression), so its multilinear extension has a closed form. The matrix $C$ is the identity, except on the wires of Section~\ref{sec:zk}'s mask, where it is zero.

Consequently the only circuit data that needs a commitment is the three local matrices. They are committed once, as a single polynomial $\idx$ of 25 variables,
\[
  \idx(x, y, T, 0) = A_T(x, y), \qquad \idx(x, y, T, 1) = A_T(x, y) + B_T(x, y),
\]
so that $\mle\idx(x, y, T, \alpha) = (A_T + \alpha B_T)(x, y)$. The commitment $C_\idx$ to $\idx$ is a constant of the verifier.

\section{The protocol}

Let $n = 22$. The prover commits to $P = (z, m_0, m_1, m_2)$, a polynomial of $n + 2$ variables: the witness $z \in \K^N$ and three uniformly random tables $m_k \in \K^N$. The transcript is seeded with $(C_\idx, y, \rootnode)$ and absorbs the commitment $C_P$ before any challenge. Write $a = Az$, $b = Bz$, $c = Cz$. The two sumchecks are the zerocheck and the lincheck of Flock~\cite{Flock26}, each with a mask.

\paragraph{1. Zerocheck.} The constraint system holds if and only if $a(x) b(x) + c(x) = 0$ on $\{0,1\}^n$. Let $g$ be the mask
\[
  g(x) = g_\varnothing + \sum_{i=1}^{n} \bigl(g_i x_i + h_i x_i^2\bigr), \qquad g_\varnothing, g_i, h_i \in \E,
\]
whose $2n + 1$ coefficients are wires of $z$ (three wires each), hence committed. The verifier sends $r \in \E^n$; the prover sends $S = \sum_x \eq(r, x) g(x)$; the verifier sends $\lambda \in \E$; and the two run the sumcheck~\cite{LFKN92} of
\[
  \sum_{x \in \{0,1\}^n} \eq(r, x) \bigl[\, \mle a(x)\, \mle b(x) + \mle c(x) + \lambda\, g(x) \,\bigr] = \lambda S .
\]
A round costs two scalars, since the factor $\eq$ is known to the verifier and the running claim fixes one coefficient. The sumcheck ends at a point $\rho \in \E^n$ with a value $e$ for the bracket. The prover sends $\hat a$, $\hat b$, $\gamma$, claimed to be $\mle a(\rho)$, $\mle b(\rho)$, $g(\rho)$, and the verifier \emph{defines} $\hat c = e + \hat a \hat b + \lambda \gamma$. Nothing is checked yet: four claims are left, all linear in $z$.

\paragraph{2. Lincheck.} Split $\rho = (\rho_\loc, \rho_{\mathrm{slot}}, \rho_{\mathrm{comp}})$ along the three parts of a wire index. The verifier sends $\alpha, \beta \in \E$ and $\pi \in \E^8$, and the four claims, with the public values, become a single one,
\begin{equation}\label{eq:lin}
  \sum_{y \in \{0,1\}^n} \comb(y)\, z(y) = v, \qquad v = \hat a + \alpha \hat b + \alpha^2 \hat c + \alpha^3 u + \alpha^4 \gamma + p ,
\end{equation}
where the weight $\comb$ is the sum of:
\begin{itemize}[nosep]
  \item $\sum_x \eq(\rho, x) (A + \alpha B + \alpha^2 C)(x, y)$, for $\hat a$, $\hat b$, $\hat c$;
  \item $\alpha^3 \eq\bigl((\rho_{\mathrm{slot}}, \rho_{\mathrm{comp}}), j\bigr)$ on the constant wire of each used block $j$, a wire that must hold 1 ($u$ is the sum of these weights);
  \item $\alpha^4$ times the weight of each mask wire in $g(\rho)$;
  \item the pins: the wires that hold $\mathtt{0x11}^{32}$, $\mathtt{0x22}^{32}$, $y$ and $\rootnode$, and the two direction bits that are fixed, each group under a power of $\beta$ and its bits under $\eq(\pi, \cdot)$, with $p$ the same combination of the public values.
\end{itemize}
The prover sends $S_k = \sum_y \comb(y) m_k(y)$ for $k = 0, 1, 2$; the verifier sends $\mu \in \E$; and the sumcheck runs on the masked table
\[
  t = z + \mu\, m_0 + \mu^2 m_1 + \mu^3 m_2, \qquad \sum_{y} \comb(y)\, t(y) = v + \mu S_0 + \mu^2 S_1 + \mu^3 S_2 ,
\]
two scalars a round. It ends at $q \in \E^n$ with a value $e' = \mle\comb(q) \cdot \mle t(q)$.

\paragraph{3. Two claims.} Split $q = (q_\loc, q_{\mathrm{slot}}, q_{\mathrm{comp}})$ as $\rho$ is. The prover sends, for each type $T$, the scalar $V_T = \mle\idx(\rho_\loc, q_\loc, T, \alpha)$. From these three values the verifier computes $\mle\comb(q)$ itself: the local part is
\[
  \eq(\rho_{\mathrm{comp}}, q_{\mathrm{comp}}) \sum_{\text{slot } s} \eq(\rho_{\mathrm{slot}}, s)\, \eq(q_{\mathrm{slot}}, s)\, V_{T(s)},
\]
and the copies, $C$, the constant wires, the pins and the mask wires are closed forms in $\rho$ and $q$. The verifier then sends $\tau \in \E^2$ and outputs the two claims
\begin{align*}
  \mle P\Bigl(q,\ \tfrac{\mu}{1 + \mu},\ \tfrac{\mu^2}{1 + \mu^2}\Bigr) &= \frac{e'}{\mle\comb(q)\,(1 + \mu)(1 + \mu^2)} && \text{on } C_P, \\
  \mle\idx(\rho_\loc, q_\loc, \tau, \alpha) &= \sum_T \eq(\tau, T)\, V_T && \text{on } C_\idx .
\end{align*}
The first is $\mle t(q)$ written as one evaluation of $P$: at those two coordinates, the four tables of $P$ are weighted $1, \mu, \mu^2, \mu^3$ up to the common factor $(1 + \mu)(1 + \mu^2)$. The verifier accepts if the precompile accepts both.

\paragraph{The proof.} $1 + 2n + 3$ scalars for the zerocheck and $3 + 2n + 3$ for the lincheck, 98 in all, and the 32-byte commitment $C_P$.

\paragraph{Why it is sound.} Each step is a standard reduction, and the masks do not weaken any of them because each is fixed before the challenge that uses it. The coefficients of $g$ are in $C_P$ and $S$ precedes $\lambda$, so a nonzero $\sum_x \eq(r,x)(ab + c)(x)$ survives for at most one $\lambda$; the claim $\gamma = g(\rho)$ is then enforced by \eqref{eq:lin} on the committed wires. The $S_k$ precede $\mu$, so a false \eqref{eq:lin} gives a nonzero polynomial of degree 3 in $\mu$. A made-up $V_T$ is caught by the claim on $C_\idx$, whose point $\tau$ is drawn after it. This note does not derive a concrete soundness bound, and takes the commitment's from~\cite{ACFY25}.

\section{Zero knowledge}\label{sec:zk}

Four things in the proof depend on the witness, and each is hidden by its own randomness.

\paragraph{Zerocheck rounds.} Round $i$ is the first whose polynomial contains $g_i$ and $h_i$, and they shift its two scalars independently. The map from the $2n + 1$ coefficients of $g$ to $S$ and the $2n$ round scalars is a bijection, so these are uniform. A sum of univariate polynomials is the mask of~\cite{Libra}; the square term is what gives a round two degrees of freedom. A multilinear mask times a public factor would not do: its round polynomial has a root the verifier knows, at which the message is the witness's own. For the bijection to be exact, the coefficients must not enter the constraint system: their wires are in no row, and $C$ is zero on them.

\paragraph{The three evaluations.} $\hat a$, $\hat b$, $\hat c$ are evaluations of the witness. Each compression carries six triples of wires $(u, v, w)$ with $u, v$ uniform in $\K$ and the row $w = u \cdot v$. They shift $(\hat a, \hat b, \hat c)$ by a term in the $u$, the $v$ and the products $uv$ that is statistically close to uniform on $\E^3$. No row of the computation reads them.

\paragraph{Lincheck rounds.} Unless $\mu \in \K$, the powers $\mu, \mu^2, \mu^3$ are a basis of $\E$ over $\K$, so $\mu m_0 + \mu^2 m_1 + \mu^3 m_2$ is a uniform table over $\E$ and so is $t$. The whole lincheck is a function of $t$, and $z$ is never evaluated: only $\mle t(q)$ reaches the commitment.

\paragraph{The opening.} Opening $C_P$ reveals Merkle leaves and folded tables, which must depend on $P$ only through the one claimed value. WHIR's first commitment interleaves $2^6$ codewords (\emph{lanes}, the blocks of the table along its top six variables), and a leaf holds every lane at one position. $P$ is therefore committed inside a table four times larger, the rest uniform over $\K$, placed for that structure:
\begin{itemize}[nosep]
  \item $P$ takes the high half of each lane it occupies. The random low half is a random polynomial of half the degree added to the lane's codeword, which makes the lane's values at any set of opened positions uniform.
  \item Every lane of $P$ is paired with a lane that is random throughout. The first fold then yields a uniform table, and everything the opening shows afterwards is a function of that table.
\end{itemize}
A claim on $P$ at a point is a claim on the larger table at the same point with $1$ and $0$ inserted at the two new coordinates.

\begin{theorem}\label{thm:zk}
Model BLAKE2s, where it hashes the leaves of $C_P$, as a random oracle. There is an efficient simulator that, given $(y, \rootnode)$ and the verifier's challenges, outputs a view such that, for every witness and over uniform challenges, the simulated and the honest views are within statistical distance $2^{-118}$ for a distinguisher that does not query the oracle, and within $2^{-118} + Q \cdot 2^{-3800}$ for one that makes $Q$ queries. The honest prover aborts with probability below $2^{-180}$, on a condition on the challenges alone.
\end{theorem}

The simulator draws the zerocheck's scalars and $(\hat a, \hat b, \hat c)$ uniformly and sets $\gamma$ from the verifier's equation; draws $t$ uniformly, derives the $S_k$ from it and from $v$, and runs the lincheck honestly on $t$; then commits to any $\K$-valued table that takes the resulting value at the resulting point, with fresh padding, and opens it honestly. The $2^{-118}$ is the probability of a set of bad challenges (for instance $\mu \in \K$), outside which the two views differ by at most $2^{-11000}$ and the oracle term. The hardest step is that the opening of the true $P$ and of the simulator's table have the same distribution: for any two $\K$-valued tables with the same claimed value there is a translation of the padding that maps one opening to the other, which is an argument on the trace form of $\E$ over $\K$. The full proof, in five lemmas, is in \code{crates/flock/membership-zk.md}.

\section{Cost}

The contract reads 2{,}384 bytes and performs 2{,}485 multiplications and 68 squarings in $\E$, and one inversion in $\K$. Its additions are XORs. Priced at 40 gas per calldata byte (the floor of~\cite{EIP7623}), and under two assumptions about a future EVM, 10 gas for a multiplication in $\E$ and 20{,}000 gas for an evaluation claim:

\begin{center}
\begin{tabular}{@{}lrr@{}}
\toprule
item & quantity & gas \\
\midrule
calldata & 2{,}384 bytes & 95{,}360 \\
multiplications and squarings & 2{,}553 & 25{,}530 \\
evaluation claims & 2 & 40{,}000 \\
\midrule
total & & 160{,}890 \\
\bottomrule
\end{tabular}
\end{center}

The table leaves out the transaction's base cost, the additions, and the hashing of the transcript. Calldata dominates, which is the reason for two scalars a round and for sending no evaluation the verifier can derive.

The prover encodes and hashes one table of $2^{26}$ words (the padded $P$), runs two sumchecks over $2^{22}$ entries, and opens that commitment at one point. The claim on $C_\idx$ concerns public data, so its opening can be produced by anyone.

\section{What is not claimed}

\begin{itemize}[nosep]
  \item \textbf{Honest verifier.} Theorem~\ref{thm:zk} is about the interactive protocol with uniform challenges. The non-interactive proof, including the proof-of-work nonces inside the opening, inherits zero knowledge by the standard argument in the programmable random oracle model, which this note does not reprove.
  \item \textbf{The commitment variant.} The argument for the opening uses that WHIR takes no out-of-domain sample on its first commitment at these parameters. A variant whose commitment includes such a sample reveals one more evaluation, which would need its own mask.
  \item \textbf{The prover's randomness} (the mask $g$, the triples, the tables $m_k$ and the padding) must be uniform and fresh for every proof.
  \item \textbf{Soundness} is argued step by step above and not quantified here.
\end{itemize}

\bibliographystyle{alpha}
\bibliography{refs}

\end{document}
